\documentclass[11pt]{article}
\usepackage[a4paper,margin=1in]{geometry}
\usepackage{amsmath,amssymb,amsthm,mathtools}
\usepackage{booktabs}
\usepackage{float}
\usepackage[hidelinks]{hyperref}
\usepackage[nameinlink,capitalise]{cleveref}
\usepackage{microtype}

\newtheorem{theorem}{Theorem}[section]
\newtheorem{proposition}[theorem]{Proposition}
\newtheorem{lemma}[theorem]{Lemma}
\newtheorem{corollary}[theorem]{Corollary}
\theoremstyle{remark}
\newtheorem{remark}[theorem]{Remark}
\crefname{lemma}{Lemma}{Lemmas}
\Crefname{lemma}{Lemma}{Lemmas}
\crefname{proposition}{Proposition}{Propositions}
\Crefname{proposition}{Proposition}{Propositions}

\newcommand{\R}{\mathbb{R}}
\newcommand{\cS}{\mathcal{S}}
\newcommand{\op}{\mathrm{op}}
\newcommand{\diag}{\operatorname{diag}}
\newcommand{\Span}{\operatorname{span}}

\title{Irreducible Channel Mixing and Exponential Calibration Laws\\
       for Matrix-Polynomial and Linear-Phase MIMO FIR Filters}
\author{Lluis Eriksson}
\date{10 August 2026}

\begin{document}
\maketitle

\begin{abstract}
Scalar Chebyshev filtering treats every vector in a block Krylov iterate with
the same polynomial, while simultaneously diagonalizable matrix coefficients
amount to independent scalar filters after a fixed channel rotation.  We study
the larger class of Hermitian matrix-polynomial filters under tangential pass
constraints.  Already at fixed channel dimension $d=3$, we construct
rationally generated signatures with quantitative full-spark margin for which
an irreducible symmetric filter has stopband leakage $O(e^{-cN})$.  In contrast,
every exactly calibrated symmetric coefficient family with any common
nontrivial invariant channel subspace has leakage at least one; this comparator
strictly contains the pairwise-commuting class and permits a fully
noncommutative $2$ by $2$ block.  A quantitative theorem covers approximately
reducible filters by charging their sampled off-block coupling together with
calibration error.  The signature margin and the admissible combined error are
both $e^{-O(N)}$, not $e^{-N^3}$.  This exponential scale is unavoidable: for
arbitrary pass nodes and signatures in the same separated bands, an explicit
scalar binomial-tail polynomial has both pass error and stopband leakage at
most $e^{-2N/81}$.  Thus constant-error separation is impossible, while the
irreducible construction achieves the correct exponential scale class.  An
affine cosine substitution gives the same robust law for fixed-latency
reciprocal linear-phase MIMO FIR filters.  We also give an exact
rational five-tap certificate with stopband norm at most $25/32$ and prove
that its advantage survives reducible calibration errors
$\delta<7/1920$.  A second exact-arithmetic certificate is calibrated from a
public triaxial pump-vibration data set: it gives leakage below $0.96$ versus
one for every exactly calibrated reducible symmetric class, while its maximum directional
error on six held-out records is $0.00385$.  In a frequency-domain-decomposition
test on the frozen held-out cospectra, filtering rotates the leading modal
direction by at most $0.222^\circ$ and changes its leading spectral ordinate by
at most $2.4\times10^{-5}$ relatively.  Numerical programs test the mechanism
and expose an ordinary graph-denoising setting in which scalar Chebyshev
filtering is instead preferable.  Tangential matrix interpolation, MIMO
filtering, scalar two-band approximation and convex FIR design are not claimed
as new; the contribution is the fixed-order reducible/irreducible separation
together with its two-sided calibration scale.
\end{abstract}

\section{Problem and contribution}

Let
\[
  P(x)=\sum_{k=0}^{N} C_k x^k, \qquad C_k\in\R^{d\times d},
\]
and let $I_-=[-1,0]$ be a stopband.  Given distinct pass points
$\lambda_j>0$ and nonzero channel signatures $v_j\in\R^d$, impose
\begin{equation}\label{eq:tangential}
  P(\lambda_j)v_j=v_j,\qquad 1\le j\le M.
\end{equation}
The leakage is
\begin{equation}\label{eq:leakage}
  \rho(P)=\sup_{x\in I_-}\|P(x)\|_{\op}.
\end{equation}
The coefficients will always be real symmetric.  They need not commute.

The principal comparator is now the larger \emph{reducible} subclass: the
coefficient matrices have a common nonzero proper invariant channel subspace.
In three dimensions, symmetry makes this equivalent to having a common
invariant line, possibly beside a fully coupled, noncommuting $2$ by $2$ block.
Pairwise-commuting coefficients form a strict subfamily: they admit a complete
common orthogonal eigenbasis and reduce to scalar filters after a fixed channel
rotation.  The question is whether genuinely irreducible channel mixing can
reduce \eqref{eq:leakage} under the same pass constraints, and at what combined
calibration/reducibility-error scale that comparison remains meaningful.  For
unit signatures define
\[
 \gamma(v_1,\ldots,v_M)
 =\min_{|J|=3}\sigma_{\min}\bigl([v_j]_{j\in J}\bigr).
\]

For a unit direction $u$ and a real-symmetric polynomial $Q$, also define its
sampled off-block coupling
\begin{equation}\label{eq:betadef}
 \beta_u(Q)=\max_j\|(I-uu^\top)Q(\lambda_j)u\|_2.
\end{equation}
It vanishes when the coefficients share the invariant line $\R u$; conversely,
because there are more pass nodes than the degree, vanishing at all pass nodes
forces that common invariant line.  Thus $\beta_u$ is a directly evaluable
relaxation of reducibility, rather than a commutator surrogate.

Our answer is an exponential calibration law.  The input signatures are full
spark and have an explicit $e^{-O(N)}$ lower margin.  This makes the reducible
obstruction quantitative under $e^{-O(N)}$ calibration and off-block errors,
while a separate scalar construction proves that no asymptotic advantage can
survive calibration errors bounded below independently of $N$.

\section{Main separation theorem}

\begin{theorem}[Robust exponential tangential separation]\label{thm:main}
There are absolute constants $C,c>0$ and $N_0\in\mathbb N$ such that for every
$N\ge N_0$, with $d=3$ and $M=N+3$, one can choose distinct rational pass
points $\lambda_1,\ldots,\lambda_M\in[1,7/2]$ and full-spark unit vectors
$v_1,\ldots,v_M\in\R^d$ satisfying
\begin{equation}\label{eq:gammafloor}
 \gamma(v_1,\ldots,v_M)
 \ge \underline\gamma_N
 :=\frac{2^{-39N}}{48(N+4)^6},
\end{equation}
with the following properties.
\begin{enumerate}
\item There is a degree-$N$ polynomial $P_N$ with real symmetric coefficient
matrices satisfying \eqref{eq:tangential} and
\[
  \rho(P_N)\le C e^{-cN}.
\]
\item If a degree-$N$ polynomial $Q_N$ has real-symmetric coefficients with a
common nonzero proper invariant subspace and satisfies the same tangential
constraints, then $\rho(Q_N)\ge1$.  If its coefficients commute pairwise, then
the stronger conclusion $Q_N(x)=I_d$ holds for every $x$.
\item More generally, let $Q_N$ be any degree-$N$ real-symmetric polynomial
obeying
\[
 \max_j\|(Q_N(\lambda_j)-I)v_j\|_2\le\delta,
\]
Then, for every unit $u\in\R^3$, with $L_N=(24e/5)^N$,
\begin{equation}\label{eq:robustmain}
 \rho(Q_N)\ge1-\frac{\sqrt3L_N}{\underline\gamma_N}
       \bigl(\delta+\beta_u(Q_N)\bigr).
\end{equation}
Consequently $\rho(Q_N)\ge3/4$ whenever
$\delta+\beta_u(Q_N)\le\underline\gamma_N/(4\sqrt3L_N)$ for some $u$.
\end{enumerate}
For every $N\ge N_0$, the coefficient matrices of $P_N$ are irreducible over
$\R$; in particular, they cannot commute pairwise.
\end{theorem}

The admissible combined error in the third item is $e^{-O(N)}$.  We split the
proof into an invariant-line obstruction, a small-leakage base construction, a
quantitatively full-spark rational perturbation, and a matching scale
obstruction.

\subsection{The invariant-line obstruction}

\begin{lemma}[Exact invariant-line obstruction]\label{lem:commuting}
Let $M\ge d+N$, let the $\lambda_j$ be distinct, and let the $v_j\in\R^d$ be
full spark.  If $Q(x)=\sum_{k=0}^{N}D_kx^k$ has real-symmetric coefficients
with a common unit eigenvector $u$ and $Q(\lambda_j)v_j=v_j$ for every $j$,
then $u^\top Q(x)u=1$ for every $x$ and $\rho(Q)\ge1$ whenever $0$ belongs to
the stopband.  If the coefficients commute pairwise, then $Q=I_d$.
\end{lemma}

\begin{proof}
Write $q(x)=u^\top Q(x)u$.  The common-eigenvector assumption and symmetry give
$Q(x)u=q(x)u$.  Taking the $u$ component of the $j$th interpolation constraint gives
\[
  (q(\lambda_j)-1)\langle u,v_j\rangle=0.
\]
A nonzero vector can be orthogonal to at most $d-1$ of the full-spark targets:
otherwise $d$ targets would lie in the $(d-1)$-dimensional hyperplane
$u^\perp$ and could not span $\R^d$.  Thus at least
$M-d+1\ge N+1$ distinct values of $\lambda_j$ are roots of $q-1$.
Consequently $q=1$, and $\|Q(0)\|_{\op}\ge|u^\top Q(0)u|=1$.

If the coefficients commute pairwise, simultaneous orthogonal diagonalization
supplies a common eigenbasis.  Applying the preceding argument to every basis
vector makes every diagonal scalar entry equal to one, hence $Q=I_d$.
\end{proof}

For completeness, $\beta_u(Q)=0$ at the $M>N$ distinct pass nodes is not merely
a sampled coincidence.  For every $z\perp u$, the scalar polynomial
$z^\top Q(x)u$ has degree at most $N$ and vanishes at all those nodes, hence is
identically zero.  Therefore every coefficient preserves $\R u$.  The
root-count implication used here and in Lemma~\ref{lem:commuting} is separately
kernel-checked in Lean/mathlib, without user axioms or admitted goals, in
\href{https://github.com/lluiseriksson/finite-sample-spectral-certificates/blob/research/endogenous-visibility/formal/SpectralCertificates/RootCount.lean}{the formal replay}.  The geometric full-spark count and the quantitative
analytic estimates remain conventional proofs, not formalized claims.

The enlargement over the old comparator is strict.  For example, coefficient
matrices of the form
\[
 D_k=a_k\oplus B_k,\qquad
 B_0=\begin{pmatrix}0&0\\0&1\end{pmatrix},\quad
 B_1=\begin{pmatrix}0&1\\1&0\end{pmatrix}
\]
share the first coordinate line but need not commute because
$[B_0,B_1]\ne0$.  Lemma~\ref{lem:commuting} still forces the scalar branch
along that line to equal one, which alone gives the leakage lower bound.

\subsection{A three-channel diagonal base filter}

Put $M=N+3$ and $S=\lceil M/3\rceil$.  Partition $M$ nodes into three
groups of sizes $s_g\in\{S-1,S\}$, $g=0,1,2$, and define
\begin{equation}\label{eq:threechannelnodes}
 \lambda_{g,r}=1+\frac{5(3r+g)}{6S},
 \qquad 0\le r<s_g.
\end{equation}
The nodes are distinct rationals in $[1,7/2]$, their global separation is at
least $5/(6S)$, and consecutive nodes within one group are separated by
$h=5/(2S)$.

Write $T_m$ for the Chebyshev polynomial of the first kind, set
$\phi(x)=2x+1$, and let $L_{g,r}$ be the Lagrange cardinal polynomial for
the nodes in group $g$.  With $m_g=N-s_g+1$, define
\begin{equation}\label{eq:multinodechebyshev}
 q_g(x)=\sum_{r=0}^{s_g-1}L_{g,r}(x)
       \frac{T_{m_g}(\phi(x))}{T_{m_g}(\phi(\lambda_{g,r}))},
 \qquad P_0(x)=\diag(q_0(x),q_1(x),q_2(x)).
\end{equation}
Each $q_g$ has degree at most $N$ and equals one at every node in its group.

\begin{lemma}\label{lem:base}
Let
\[
 \eta_0=\operatorname{arcosh}(3),\quad
 \eta_1=\operatorname{arcosh}(8),\quad
 \Delta\eta=\eta_1-\eta_0,
 \quad K=\frac{36e}{5},
\]
and set
\[
 c=\frac{2\eta_0-\log K}{3}>0,\qquad
 C_0=2\exp\!\left(\frac23(\log K+\eta_0)\right).
\]
For every $N\ge7$,
\[
 \sup_{x\in I_-}\|P_0(x)\|_{\op}\le C_0e^{-cN}.
\]
Moreover, at every pass node,
\[
 \|P_0(\lambda_j)\|_{\op}
 \le B_N:=2(4e)^{S-1}e^{N\Delta\eta}.
\]
\end{lemma}

\begin{proof}
For $x\in I_-$, $|x-\lambda_{g,r}|\le9/2$.  Equispacing within a group gives
\[
 \sum_{r=0}^{s_g-1}|L_{g,r}(x)|
 \le \frac{(18S/5)^{s_g-1}}{(s_g-1)!}.
\]
Here we used
$\sum_r[r!(s_g-1-r)!]^{-1}=2^{s_g-1}/(s_g-1)!$.
For $N\ge7$, $S\ge4$, $s_g\ge S-1$, and hence
$S/(s_g-1)\le2$.  The bound $n!\ge(n/e)^n$ therefore gives
\begin{equation}\label{eq:stoplagrange}
 \sum_r|L_{g,r}(x)|\le K^{s_g-1}.
\end{equation}
Also, for $\lambda\ge1$,
\[
 T_m(\phi(\lambda))\ge\tfrac12e^{m\eta_0}.
\]
Combining this with \eqref{eq:stoplagrange} yields
\[
 |q_g(x)|\le2K^{s_g-1}e^{-(N-s_g+1)\eta_0}.
\]
Since $s_g-1\le S-1\le(N+2)/3$, the claimed $C_0e^{-cN}$ follows.
The constant $c$ is positive because $e<3$ gives
$K<108/5<e^{2\eta_0}=(3+\sqrt8)^2$.

At another pass node the numerator distances in the Lagrange factors are at
most $5/2$, so the same calculation gives
$\sum_r|L_{g,r}(\lambda_j)|\le(4e)^{s_g-1}$.
The Chebyshev ratio is at most $2e^{N\Delta\eta}$, proving the pass-node
bound.  The operator norm of a diagonal polynomial is its largest entry.
\end{proof}

\subsection{An explicit full-spark perturbation}

Index all pass nodes by $j=1,\ldots,M$, let $g(j)$ be the coordinate group of
node $j$, and set $t_j=j/(M+1)$.  Define
\begin{equation}\label{eq:momentperturb}
 w_j=(1,t_j,t_j^2)^\top,\quad
 \varepsilon_N=2^{-13N},\quad
 \widetilde v_j=e_{g(j)}+\varepsilon_Nw_j,\quad
 v_j=\frac{\widetilde v_j}{\|\widetilde v_j\|_2}.
\end{equation}

\begin{lemma}[Quantitative rational full spark]\label{lem:fullspark}
The vectors in \eqref{eq:momentperturb} are full spark and satisfy
\eqref{eq:gammafloor} for every $N\ge7$.
\end{lemma}

\begin{proof}
Fix any three indices and write
\[
 D(z)=\det(e_{g(j_1)}+zw_{j_1},e_{g(j_2)}+zw_{j_2},
                 e_{g(j_3)}+zw_{j_3})=\sum_{r=0}^3d_rz^r.
\]
The cubic coefficient is the nonzero Vandermonde determinant of the three
$w_j$.  Every $d_r$ is rational with denominator dividing $(M+1)^6$;
therefore every nonzero coefficient has magnitude at least $(M+1)^{-6}$.
Expanding the determinant into its six signed products also gives
$|d_r|\le18$.

Let $r_0$ be the first nonzero coefficient.  For $N\ge7$,
$\varepsilon_N\le[108(M+1)^6]^{-1}$, and hence
\[
 \sum_{r>r_0}|d_r|\varepsilon_N^{r-r_0}
 \le18(\varepsilon_N+\varepsilon_N^2+\varepsilon_N^3)
 \le\frac1{2(M+1)^6}.
\]
It follows that
\[
 |D(\varepsilon_N)|
 \ge\frac{\varepsilon_N^{r_0}}{2(M+1)^6}
 \ge\frac{\varepsilon_N^3}{2(M+1)^6}.
\]
Each unnormalized column has norm below two, so normalization loses at most a
factor eight in the determinant.  For a matrix with three unit columns the
product of its two largest singular values is at most three.  Its smallest
singular value is consequently at least
\[
 \frac{\varepsilon_N^3}{48(M+1)^6}
 =\frac{2^{-39N}}{48(N+4)^6}.
\]
The elementary inequality on $\varepsilon_N$ holds at $N=7$ and its left-to-
right ratio decreases thereafter.
\end{proof}

\begin{lemma}[Stable symmetric interpolation]\label{lem:stability}
For every $N\ge7$, there is a real-symmetric coefficient polynomial $P_N$
satisfying $P_N(\lambda_j)\widetilde v_j=\widetilde v_j$ and
\[
 \sup_{x\in I_-}\|P_N(x)-P_0(x)\|_{\op}
 \le e^{-cN}.
\]
\end{lemma}

\begin{proof}
Let $\cS_N$ be the finite-dimensional space of degree-at-most-$N$ polynomials
with real symmetric coefficients, and define
\[
 L_\varepsilon:\cS_N\longrightarrow(\R^d)^M,
 \qquad L_\varepsilon(P)
 =\bigl(P(\lambda_j)(e_{g(j)}+\varepsilon w_j)\bigr)_{j=1}^M.
\]
At $\varepsilon=0$, this map is onto.  Indeed, for an off-diagonal entry
$p_{rs}=p_{sr}$, arbitrary output data prescribe values at the union of nodes
assigned to coordinates $r$ and $s$.  That union contains at most
$2S\le N+1$ nodes for $N\ge7$.
Ordinary degree-$N$ Lagrange interpolation is therefore onto on those nodes.
A diagonal entry sees at most $S$ nodes.

Equip the output space with the maximum Euclidean block norm and $\cS_N$ with
\[
 \|P\|_X=\max\left\{\sup_{x\in I_-}\|P(x)\|_{\op},
                     \max_j\|P(\lambda_j)\|_{\op}\right\}.
\]
This also gives an explicit linear right inverse.  To record its norm, let
$y=(y_j)_{j=1}^M$ satisfy $\max_j\|y_j\|_2\le1$.  For $r\ne g$, define the
entry $(R_Ny)_{rg}=(R_Ny)_{gr}$ by scalar Lagrange interpolation on the union
of groups $r$ and $g$: its value is $(y_j)_r$ at a node in group $g$ and
$(y_j)_g$ at a node in group $r$.  A diagonal entry uses only its own group.
These prescriptions are consistent because the groups are disjoint, and the
resulting symmetric polynomial obeys $L_0R_Ny=y$.

For any such scalar problem, write its $t\le2S$ nodes in increasing order.
Their separation $\delta=5/(6S)$ implies that the denominator of the $i$th
cardinal polynomial is at least
\[
 \delta^{t-1}i!(t-1-i)!.
\]
On $I_-$ and at every pass node, all numerator distances are at most $9/2$.
Put $n=t-1$.  Since every scalar problem contains at least $S-1$ nodes,
$S/n\le2$ for $S\ge4$.  Consequently the sum of the absolute cardinal
functions is bounded by
\[
 \left(\frac{27S}{5}\right)^{t-1}
 \sum_{i=0}^{t-1}\frac1{i!(t-1-i)!}
 =\left(\frac{27S}{5}\right)^{t-1}
   \frac{2^{t-1}}{(t-1)!}
 \le\left(\frac{108e}{5}\right)^n.
\]
The operator norm of a $3$ by $3$ matrix is at most three times its largest
entry.  Hence, with the norm $\|\cdot\|_X$ above,
\begin{equation}\label{eq:rightinversebound}
 \|R_N\|_{Y\to X}
 \le6\left(\frac{108e}{5}\right)^{2S-1}
 =:\mathcal R_N.
\end{equation}
Since $S=\lceil(N+3)/3\rceil$, this explicit envelope satisfies
$\mathcal R_N=\exp(O(N))$.  Retaining the factorial is the step that improves
the old superexponential perturbation scale.

Write $L_\varepsilon=L_0+\varepsilon K_N$.  Since
$\|w_j\|_2\le\sqrt3$, $\|K_N\|_{X\to Y}\le\sqrt3$.  Hence
\[
 L_\varepsilon R_N=I+E_N,\qquad
 \|E_N\|\le\sqrt3\,\varepsilon\mathcal R_N.
\]
For $\varepsilon=\varepsilon_N$ and $N\ge7$, this is below $1/2$.  Indeed,
using $e<11/4$, $e^{\Delta\eta}<4$, $S\le(N+5)/3$, and
$108e/5<64$ gives
\begin{equation}\label{eq:powerenvelopes}
 \mathcal R_N\le2^{4N+17},\qquad
 B_N\le2^{(10N+11)/3}.
\end{equation}
Thus $\|E_N\|\le2^{18-9N}<1/2$, so $(I+E_N)^{-1}$ exists and has norm at
most two.

Since $L_0(P_0)=(e_{g(j)})_j$, the residual at the perturbed targets is
\[
 r_N=(\widetilde v_j)_j-L_{\varepsilon_N}(P_0)
 =\varepsilon_N\bigl((I-P_0(\lambda_j))w_j\bigr)_j.
\]
By Lemma~\ref{lem:base},
\[
 \|r_N\|_Y\le\varepsilon_N\sqrt3(1+B_N).
\]
Set
\[
 \Delta P_N=R_N(I+E_N)^{-1}r_N,
 \qquad P_N=P_0+\Delta P_N.
\]
Then $L_{\varepsilon_N}(P_N)=(\widetilde v_j)_j$ and
\[
 \|\Delta P_N\|_X
 \le2\sqrt3\,\mathcal R_N\varepsilon_N(1+B_N)\le e^{-cN}.
\]
Indeed, $1+B_N\le2B_N$, so \eqref{eq:powerenvelopes} bounds the correction
by $2^{(71-17N)/3}$.  Since $c<\log2$, the ratio of this envelope to
$e^{-cN}$ is at most $2^{(71-14N)/3}\le2^{-9}$ for $N\ge7$.
Every coefficient remains real symmetric because $R_N$ maps into $\cS_N$.
\end{proof}

\begin{lemma}[Quantitative near-reducibility obstruction]\label{lem:robustcommuting}
For the nodes and signatures above, let $Q$ be any degree-$N$ polynomial with
real-symmetric coefficients and suppose
\[
 \max_j\|(Q(\lambda_j)-I)v_j\|_2\le\delta.
\]
Then \eqref{eq:robustmain} holds for every unit $u\in\R^3$.
\end{lemma}

\begin{proof}
Fix a unit vector $u$ and put $q(x)=u^\top Q(x)u$.  At each pass node set
$r_j=(I-uu^\top)Q(\lambda_j)u$.  Symmetry and the orthogonal decomposition of
$v_j$ give
\[
 u^\top(Q(\lambda_j)-I)v_j
 =(q(\lambda_j)-1)u^\top v_j+r_j^\top(I-uu^\top)v_j.
\]
The left-hand side has magnitude at most $\delta$, while
$\|r_j\|_2\le\beta_u(Q)$.  If three signatures all had
$|u^\top v_j|<\underline\gamma_N/\sqrt3$, then the Euclidean norm of the
three overlaps would be below $\underline\gamma_N$, contradicting
\eqref{eq:gammafloor}.  Thus at most two signatures have such a small overlap,
and at $N+1$ distinct nodes
\[
 |q(\lambda_j)-1|
 \le\frac{\sqrt3(\delta+\beta_u(Q))}{\underline\gamma_N}.
\]
For any selected $N+1$ nodes, ordered increasingly, their separation is at
least $5/(6S)$.  The sum of the absolute degree-$N$ Lagrange cardinal values
at zero is at most
\[
 \left(\frac{42S}{5}\right)^N\frac1{N!}
 \le\left(\frac{24e}{5}\right)^N=L_N,
\]
where $S/N\le4/7$ for $N\ge7$.  Hence
$|q(0)-1|\le\sqrt3L_N(\delta+\beta_u(Q))/\underline\gamma_N$.
Finally $\|Q(0)\|_{\op}\ge|u^\top Q(0)u|=|q(0)|$, and $0\in I_-$.
\end{proof}

\begin{proof}[Proof of \cref{thm:main}]
The normalized and unnormalized interpolation constraints are equivalent.
Combine Lemmas~\ref{lem:base}, \ref{lem:fullspark}, and \ref{lem:stability}.
The first conclusion holds with $C=C_0+1$ and the $c>0$ of
Lemma~\ref{lem:base}.  For the second conclusion, a common nontrivial invariant
subspace of symmetric $3$ by $3$ matrices yields a common invariant line: use
the subspace itself if it is one-dimensional and its invariant orthogonal
complement otherwise.  Lemma~\ref{lem:commuting} then gives leakage at least
one, and gives the identity conclusion in the commuting subcase.
The third is Lemma~\ref{lem:robustcommuting}.  Choose $N_0\ge7$ so that
$Ce^{-cN_0}<3/4$.  Then the constructed coefficients cannot commute pairwise
or have any common nontrivial invariant subspace for $N\ge N_0$, and the
robust comparison is strict at the displayed combined-error radius.
\end{proof}

\subsection{Why the exponential calibration scale is unavoidable}

The preceding radius is small, but its exponential order cannot be replaced by
a constant.  The obstruction is already scalar and is independent of the
signature geometry.

\begin{proposition}[Universal scalar noise barrier]\label{prop:scalarbarrier}
For every $N\ge1$, arbitrary pass nodes $\lambda_j\in[1,7/2]$, and arbitrary
unit signatures $v_j$, there is a scalar degree-$N$ polynomial $s_N$ such that
\[
 \sup_{x\in[-1,0]}|s_N(x)|\le e^{-2N/81},\qquad
 \max_j|s_N(\lambda_j)-1|\le e^{-2N/81}.
\]
Consequently the commuting symmetric polynomial $s_NI$ has both leakage and
maximum tangential calibration residual at most $e^{-2N/81}$.
\end{proposition}

\begin{proof}
Put $r(x)=2(x+1)/9$ and $k_N=\lceil N/3\rceil$.  Define the binomial-tail
polynomial
\[
 s_N(x)=\sum_{k=k_N}^N\binom{N}{k}r(x)^k(1-r(x))^{N-k}.
\]
On the stopband $0\le r\le2/9$, while on the pass interval
$4/9\le r\le1$.  Hoeffding's inequality for a binomial random variable
\cite{Hoeffding1963} gives
\[
 \Pr\{X/N\ge1/3\}\le e^{-2N/81}\quad(r\le2/9)
\]
and
\[
 \Pr\{X/N<1/3\}\le e^{-2N/81}\quad(r\ge4/9).
\]
The ceiling in $k_N$ only strengthens both deviations.  These are precisely
the two claimed bounds.
\end{proof}

Thus exact feasibility yields a class separation, exponentially accurate
feasibility retains it for the constructed signatures, and exponentially
accurate scalar filters eventually erase it for every signature family.  The
sufficient and impossible tolerances do not have matching constants, but both
have $-\log\delta=\Theta(N)$; an $N$-independent calibration floor is ruled out.

\section{Block-Krylov meaning}

The theorem has a direct interpretation for block polynomial filtering.  Let
$H=\sum_j\lambda_jq_jq_j^\top$ be real symmetric and let $B\in\R^{n\times d}$
be a starting block.  For $P(x)=\sum_{k=0}^NC_kx^k$, define the right-coefficient
block filter
\begin{equation}\label{eq:blockfilter}
 \mathfrak F_P(H,B)=\sum_{k=0}^N H^kBC_k.
\end{equation}
Its $j$th spectral row is
\begin{equation}\label{eq:spectralrow}
 q_j^\top\mathfrak F_P(H,B)=(q_j^\top B)P(\lambda_j).
\end{equation}

\begin{proposition}[Pass preservation and stop contraction]
\label{prop:krylov}
Let $\mathcal T$ and $\mathcal S$ index target and stop eigenpairs.  Assume the
rows $b_j^\top=q_j^\top B$ obey
\[
 b_j^\top P(\lambda_j)=b_j^\top \quad(j\in\mathcal T),
 \qquad \lambda_j\in I_- \quad(j\in\mathcal S).
\]
Then the target spectral component is preserved exactly, and
\[
 \|Q_\mathcal S^\top\mathfrak F_P(H,B)\|_F
 \le \rho(P)\|Q_\mathcal S^\top B\|_F,
\]
where $Q_\mathcal S$ has columns $q_j$, $j\in\mathcal S$.
If the $C_k$ are pairwise-commuting symmetric matrices, one fixed orthogonal
channel rotation reduces \eqref{eq:blockfilter} to $d$ independent scalar
polynomial filters.
\end{proposition}

\begin{proof}
The pass statement follows from \eqref{eq:spectralrow}.  For the stop part,
\[
 \|Q_\mathcal S^\top\mathfrak F_P(H,B)\|_F^2
 =\sum_{j\in\mathcal S}\|b_j^\top P(\lambda_j)\|_2^2
 \le\rho(P)^2\sum_{j\in\mathcal S}\|b_j\|_2^2.
\]
For the final claim, simultaneously diagonalize all $C_k=U\Lambda_kU^\top$.
Then the columns of $BU$ are acted on by the scalar polynomials whose
coefficients are the diagonal entries of the $\Lambda_k$.
\end{proof}

Because every coefficient in \cref{thm:main} is symmetric, its column
constraints transpose to the row constraints in Proposition~\ref{prop:krylov}.  The
theorem therefore gives block spectral instances with exponentially small stop
contraction that no fixed-basis collection of scalar pass-preserving filters
can reproduce.  This is an approximation-theoretic separation; an end-to-end
runtime advantage also depends on filter construction, recurrence stability,
orthogonalization, and the cost model.

\section{Reciprocal linear-phase MIMO FIR meaning}

The polynomial separation is equivalently a fixed-order statement for a
standard multichannel FIR class.  Let
\[
 x(y)=\frac{9y+5}{4},\qquad G_N(\omega)=P_N(x(\cos\omega)).
\]
Because $G_N$ is a real symmetric matrix cosine polynomial of degree $N$, it
has a unique expansion
\[
 G_N(\omega)=A_N+2\sum_{r=1}^{N}A_{N-r}\cos(r\omega),
 \qquad A_r=A_r^\top.
\]
The causal response
\begin{equation}\label{eq:firresponse}
 H_N(e^{i\omega})=\sum_{r=0}^{2N}A_re^{-ir\omega}
                  =e^{-iN\omega}G_N(\omega),
 \qquad A_r=A_{2N-r},
\end{equation}
therefore has $2N+1$ real symmetric palindromic taps and common group delay
$N$.

\begin{corollary}[Fixed-latency MIMO FIR separation]\label{cor:fir}
For every $N\ge N_0$ there are $M=N+3$ distinct frequencies
$\omega_j\in[0,\arccos(-1/9)]$ and full-spark directions
$v_j\in\R^{3}$ for which a causal response of the form
\eqref{eq:firresponse} satisfies
\[
 H_N(e^{i\omega_j})v_j=e^{-iN\omega_j}v_j,
 \qquad
 \sup_{\omega\in[\arccos(-5/9),\pi]}
       \|H_N(e^{i\omega})\|_{\op}\le Ce^{-cN}.
\]
Every response with the same real-symmetric palindromic tap structure whose
taps have a common nontrivial invariant channel subspace and which satisfies
the same calibrations has high-frequency norm at least one.  If the taps
commute pairwise, it is the pure delay $e^{-iN\omega}I$.
More generally, put $G(\omega)=e^{iN\omega}H(e^{i\omega})$ and
\[
 \beta_u^{\rm FIR}(H)=
 \max_j\|(I-uu^\top)G(\omega_j)u\|_2.
\]
For any symmetric palindromic comparator with phase-corrected calibration
residual at most $\delta$, its high-frequency norm is at least the right-hand
side of \eqref{eq:robustmain} with $\beta_u(Q_N)$ replaced by
$\beta_u^{\rm FIR}(H)$.  Conversely, for arbitrary calibration directions and
frequencies in the same pass range, there is a scalar palindromic response with
both phase-corrected pass error and high-frequency norm at most $e^{-2N/81}$.
\end{corollary}

\begin{proof}
Take the data of \cref{thm:main} and choose
$\cos\omega_j=(4\lambda_j-5)/9$.  The pass interval $[1,7/2]$ maps into
$[-1/9,1]$, while the stop interval $[-1,0]$ maps onto $[-1,-5/9]$.
Equation~\eqref{eq:firresponse} gives the calibration and transfers the norm
bound unchanged.  The affine substitution is invertible on coefficient
spaces, as is the change between powers of $\cos\omega$ and the cosine basis.
Thus the transformed coefficient matrices commute pairwise exactly when the
original polynomial coefficients do, and the invertible coefficient changes
preserve common invariant subspaces.  Palindromy makes the FIR taps an
invertible relabelling and rescaling of the cosine coefficients.  Applying
Lemma~\ref{lem:commuting} therefore forces a unit-gain scalar branch for every
reducible response and the full identity in the commuting subcase.
Lemma~\ref{lem:robustcommuting} transfers both residual and off-block bounds, and
Proposition~\ref{prop:scalarbarrier} transfers the universal scalar response.
\end{proof}

The reducible comparator in Corollary~\ref{cor:fir} is strictly broader than a
bank of scalar linear-phase filters in one fixed basis: it also allows an
arbitrary coupled, noncommuting $2$ by $2$ symmetric block beside the protected
scalar branch.  The result does not compare implementation cost, quantization
robustness, rational/state-space filters, or nonsymmetric MIMO filters, and the
asymptotic calibrated directions are constructed rather than taken from a
deployed system.

\section{Exact finite witnesses}

\subsection{A full-spark quadratic certificate}

The asymptotic mechanism already occurs with rational data at degree two.  Take
\[
 (\lambda_1,\lambda_2,\lambda_3,\lambda_4)
 =\left(\tfrac12,1,\tfrac32,\tfrac72\right),
 \qquad
 (v_1,v_2,v_3,v_4)
 =\left(
 \binom{1}{2},\binom{0}{1},\binom{-1}{2},\binom{-1}{1}
 \right).
\]
Every pair of targets is linearly independent, hence the tuple is full spark in
$\R^2$.  Set $P(x)=C_0+xC_1+x^2C_2$, with
\[
 C_0=\begin{pmatrix}61/96&7/48\\7/48&61/96\end{pmatrix},\quad
 C_1=\begin{pmatrix}7/24&7/48\\7/48&77/96\end{pmatrix},\quad
 C_2=\begin{pmatrix}-7/24&-7/24\\-7/24&-7/16\end{pmatrix}.
\]
Direct substitution verifies $P(\lambda_j)v_j=v_j$ for all four targets, and
\[
 [C_0,C_1]
 =\begin{pmatrix}0&343/4608\\-343/4608&0\end{pmatrix}\ne0.
\]

\begin{proposition}[Exact quadratic gap]\label{prop:quadratic}
The polynomial above obeys
\[
 \sup_{x\in[-1,0]}\|P(x)\|_{\op}\le\frac{25}{32}<1.
\]
Every pairwise-commuting real-symmetric quadratic polynomial satisfying the
same four tangential constraints is identically $I_2$ and has leakage one.
\end{proposition}

\begin{proof}
For the upper bound, put $x=t-1$ and express the leading principal minor and
determinant of $(25/32)I_2\pm P(t-1)$ in the Bernstein basis on $t\in[0,1]$.
The leading-minor coefficients are
\[
 \left(\tfrac{35}{48},\tfrac7{24},\tfrac7{48}\right),\qquad
 \left(\tfrac56,\tfrac{61}{48},\tfrac{17}{12}\right),
\]
and the determinant coefficients are, respectively,
\begin{align*}
 &\left(\tfrac{1421}{1536},\tfrac{2597}{6144},\tfrac{4655}{27648},
          \tfrac{931}{18432},0\right),\\
 &\left(\tfrac1{16},\tfrac{1711}{3072},\tfrac{3835}{3456},
          \tfrac{3707}{2304},\tfrac{1525}{768}\right).
\end{align*}
All are nonnegative and both leading minors are positive.  Sylvester's
criterion therefore gives $(25/32)I_2\pm P(x)\succeq0$ throughout the
stopband, proving the norm bound.  The zero final coefficient in the first
determinant shows that the bound is attained at an endpoint rather than being
a grid artifact.

For the commuting class, apply Lemma~\ref{lem:commuting} with
$M=4=d+N$.  Every scalar common-eigenbasis entry has three distinct roots of
$q-1$ and is identically one.
\end{proof}

The gap is stable under imperfect calibration.  Write
$\widehat v_j=v_j/\|v_j\|_2$.

\begin{proposition}[Robust near-reducibility obstruction]\label{prop:robustquadratic}
Let $Q$ be any real-symmetric quadratic polynomial satisfying
\[
 \|Q(\lambda_j)\widehat v_j-\widehat v_j\|_2\le\delta,
 \qquad j=1,\ldots,4.
\]
For every unit $u\in\R^2$, define $\beta_u(Q)$ by \eqref{eq:betadef}.  Then
\[
 \sup_{x\in[-1,0]}\|Q(x)\|_{\op}
 \ge 1-60\bigl(\delta+\beta_u(Q)\bigr).
\]
In particular, a reducible comparator has $\beta_u(Q)=0$ for a common
invariant direction and, if $\delta<7/1920$, its leakage is strictly larger
than the certified $25/32$ leakage of the irreducible witness.
\end{proposition}

\begin{proof}
For any two distinct normalized targets $a,b$ in the tuple,
$|a^\top b|^2\le9/10$.  Hence their two-column smallest singular value obeys
\[
 \sigma_{\min}^2=1-|a^\top b|
 \ge\frac{1-|a^\top b|^2}{2}\ge\frac1{20}>\frac1{25}.
\]
For every unit vector $u$, at least one member of any pair therefore has
$|u^\top\widehat v_j|\ge1/(5\sqrt2)$.  Thus at most one target has smaller
overlap.
Fix $u$ and put $q(x)=u^\top Q(x)u$.  The same symmetric off-block
decomposition used in Lemma~\ref{lem:robustcommuting} shows that at three
distinct pass nodes
\[
 |q(\lambda_j)-1|\le5\sqrt2\bigl(\delta+\beta_u(Q)\bigr).
\]
For each of the four possible three-node subsets of
$\{1/2,1,3/2,7/2\}$, the sum of the absolute Lagrange weights at $x=0$ is
at most eight.  Hence
$|q(0)-1|<60(\delta+\beta_u(Q))$, so
$\|Q(0)\|_{\op}\ge1-60(\delta+\beta_u(Q))$.  Finally,
$1-25/32=7/32$ gives the stated threshold.
\end{proof}

The exact-arithmetic replay
\texttt{verification/verify\_quadratic\_filter\_witness.py} checks all
constraints, target minors, commutators, Bernstein coefficients, and the
rational robust-calibration constants without a numerical solver.

\subsection{An exact five-tap MIMO FIR certificate}

The quadratic witness has an exact causal FIR realization.  Substitute
$x=(9\cos\omega+5)/4$ and define the palindromic taps
\[
 B_0=B_4=\begin{pmatrix}-189/512&-189/512\\-189/512&-567/1024\end{pmatrix},
\]
\[
 B_1=B_3=\begin{pmatrix}-63/128&-21/32\\-21/32&-21/64\end{pmatrix},\quad
 B_2=\begin{pmatrix}-149/768&-665/768\\-665/768&-235/1536\end{pmatrix}.
\]
Their response obeys the exact identity
\[
 \sum_{r=0}^{4}B_re^{-ir\omega}
 =e^{-2i\omega}\bigl(B_2+2B_1\cos\omega+2B_0\cos2\omega\bigr)
 =e^{-2i\omega}P\!\left(\frac{9\cos\omega+5}{4}\right).
\]
Consequently it preserves the four directions, up to their common delay
phase, at cosine values $-1/3,-1/9,1/9,1$ and has operator norm at most $25/32$
on $\omega\in[\arccos(-5/9),\pi]$.  Every reducible symmetric palindromic
five-tap response meeting the same four calibrations is the pure two-sample
delay and has norm one.  Moreover $[B_0,B_1]\ne0$.  All statements, including
the tap conversion, are checked by the same exact-arithmetic replay.  The same
change of variables transfers Proposition~\ref{prop:robustquadratic}: if a
five-tap comparator obeys
$\|e^{2i\omega_j}H(e^{i\omega_j})\widehat v_j-\widehat v_j\|_2\le\delta$,
its high-frequency norm is at least
$1-60(\delta+\beta_u^{\rm FIR}(H))$ for every unit $u$.

\subsection{An affine two-constraint precursor}

For degree one and $d=2$, take pass data
\[
 \lambda_1=\tfrac12,\quad v_1=(1,0)^\top,
 \qquad \lambda_2=1,\quad v_2=(1,1)^\top,
\]
and $P(x)=A+xB$ with
\[
 A=\begin{pmatrix}4/5&1/5\\1/5&3/10\end{pmatrix},\qquad
 B=\begin{pmatrix}2/5&-2/5\\-2/5&9/10\end{pmatrix}.
\]
The constraints hold exactly and
\[
 [A,B]=\begin{pmatrix}0&-1/10\\1/10&0\end{pmatrix}.
\]
Convexity of $x\mapsto\|A+xB\|_{\op}$ puts its maximum over $I_-$ at an
endpoint.  Direct diagonalization gives
\[
 \rho(P)=\max\{\|A\|_{\op},\|A-B\|_{\op}\}
 =\frac{1+\sqrt{61}}{10}<0.882.
\]
If $A$ and $B$ commute, their common eigenbasis contains a direction with
nonzero overlap with both nonorthogonal, nonparallel targets.  Its scalar
affine entry equals one at two points and hence everywhere.  The commuting
leakage is at least one.

\section{Numerical stress test}

The repository script \texttt{research/route\_c\_scaling\_pilot.py} solves a
grid SDP with symmetric coefficient variables.  For degree $3$, dimension $9$,
$12$ perturbed targets and seed 7, the results are:
\begin{center}
\begin{tabular}{@{}rrrr@{}}
\toprule
Perturbation & SDP objective & Validation maximum & Commutator norm\\
\midrule
$0.001$ & $0.2983490$ & $0.2983682$ & $0.466925$\\
$0.010$ & $0.3199793$ & $0.3200281$ & $0.391179$\\
$0.030$ & $0.3977121$ & $0.3977588$ & $0.708604$\\
\bottomrule
\end{tabular}
\end{center}
The optimization uses 81 stopband points and validation uses 4001 independent
points.  Clarabel labels these solves \texttt{optimal\_inaccurate}; they are
diagnostics, not proof.  The exact witnesses and \cref{thm:main} do not depend
on solver output.

\subsection{Controlled block-iteration benchmark}

The script \texttt{research/route\_c\_block\_krylov\_benchmark.py} embeds the
quadratic certificate in a $256$-dimensional Hermitian eigenproblem.  The four
pass eigenvalues are $1/2,1,3/2,2$, the other $252$ eigenvalues fill $[-1,0]$,
and the initial stop/pass Frobenius ratio is five.  We repeatedly apply the
degree-two matrix filter and compare it with

\[
 q(x)=\frac{8x^2+8x+1}{7},
\]

the scalar degree-two Chebyshev stopband filter normalized only at $x=1/2$.
Let $\sin\theta$ denote the largest principal-angle sine from the filtered
block to the clean two-dimensional block defined by the four prescribed pass
rows.  The scalar pass error below is minimized over one global rescaling, so
it does not merely report scalar amplitude growth.

\begin{center}
\begin{tabular}{@{}rrrrr@{}}
\toprule
Iterations & Matrix stop/pass & Matrix $\sin\theta$ & Scalar pass error & Scalar $\sin\theta$\\
\midrule
$0$  & $5.0000$ & $0.9918$ & $0.0000$ & $0.9918$\\
$1$  & $2.2207$ & $0.9638$ & $0.7518$ & $0.7847$\\
$5$  & $0.3752$ & $0.5622$ & $0.9195$ & $0.9802$\\
$10$ & $0.0760$ & $0.1385$ & $0.9199$ & $0.9819$\\
$20$ & $0.00434$ & $0.00810$ & $0.9199$ & $0.9820$\\
\bottomrule
\end{tabular}
\end{center}

The matrix pass error is below $1.2\times10^{-10}$ through ten iterations and
is $7.8\times10^{-4}$ at iteration 20 because floating-point errors are
amplified outside the preserved directions; exact arithmetic preserves them.
The scalar filter has stopband maximum $1/7$ and initially suppresses
contamination much faster, but its pass multipliers are
$(1,17/7,31/7,127/7)$; iteration consequently rotates the clean block away
from the prescribed one.  The fair commuting comparator under all four exact
tangential constraints is the identity and remains at stop/pass ratio $5$ and
$\sin\theta=0.9918$.  This controlled construction demonstrates both the
separation and its recurrence-conditioning limitation, not a runtime advantage
on a natural problem.

\subsection{A measured triaxial calibration and exact continuum certificate}

We next use VBL-VA001, an open lab-scale pump-vibration data set with triaxial
acceleration records sampled at $20$ kHz \cite{AtmajaEtAl2024,VBLVA001}.
The task is frequency-domain-decomposition (FDD) preprocessing, not fault
classification.  FDD extracts mode-shape estimates as leading singular vectors
of cross-spectral matrices at spectral peaks \cite{BrinckerZhangAndersen2001};
filter-induced perturbations are therefore a scientifically relevant failure
mode, and filtering artifacts can corrupt modal estimates \cite{WynneEtAl2022}.
We selected twelve
normal-condition records by the archive prefixes \texttt{normal\_000} and
\texttt{normal\_001}, using all six records under the first prefix for fitting
and retaining all six under the second prefix.  The file names, CRC-32 and
SHA-256 digests were frozen before optimization.

Each record was polyphase-decimated by eight, then divided into length-$1024$
Hann windows with $50\%$ overlap.  We averaged the real $3$ by $3$ cospectral
matrix and, using the training split only, selected the five largest-trace
peaks between $20$ and $1000$ Hz subject to an eight-bin separation.  Their
leading cospectral eigenvectors are the channel signatures.  The frequency
node is $x=\cos(2\pi f/2500)$.  Table~\ref{tab:vblcalibration} reports the
six-decimal rationalization used by the exact certificate; every displayed
component and node differs from its measured value by less than
$5\times10^{-7}$.

\begin{lemma}[Modal-component invariance]\label{lem:fddinvariance}
Let the cospectral matrix at a modal frequency have the decomposition
$G_j=\lambda_jv_jv_j^\top+R_j$, where $\lambda_j\ge0$, $\|v_j\|_2=1$, and
$R_j\succeq0$.  If a real-symmetric response $P$ obeys
$P(x_j)v_j=v_j$, then the reciprocal palindromic FIR response
$H(e^{\mathrm{i}\omega})=e^{-\mathrm{i}N\omega}P(\cos\omega)$ obeys
\[
 H(e^{\mathrm{i}\omega_j})G_jH(e^{\mathrm{i}\omega_j})^*
 =\lambda_jv_jv_j^\top+P(x_j)R_jP(x_j).
\]
Thus the rank-one modal component, its channel shape, and its amplitude are
preserved exactly, up to one common $N$-sample delay, while the residual
cospectrum is filtered.
\end{lemma}

\begin{proof}
The scalar phase cancels in $H G_j H^*$, and symmetry gives
$P(x_j)v_jv_j^\top P(x_j)=v_jv_j^\top$.
\end{proof}

\begin{table}[H]
\centering
\small
\begin{tabular}{@{}rrrrr@{}}
\toprule
$f$ (Hz) & $x$ & $v_x$ & $v_y$ & $v_z$\\
\midrule
$100.0977$ & $ 0.968522$ & $-0.054531$ & $ 0.996129$ & $-0.068944$\\
$200.1953$ & $ 0.876070$ & $ 0.980644$ & $ 0.167125$ & $ 0.102012$\\
$744.6289$ & $-0.296151$ & $ 0.152303$ & $ 0.986844$ & $-0.054250$\\
$844.7266$ & $-0.524590$ & $ 0.869001$ & $-0.350335$ & $ 0.349432$\\
$944.8242$ & $-0.720003$ & $ 0.606651$ & $ 0.716228$ & $-0.344951$\\
\bottomrule
\end{tabular}
\caption{Training frequencies and rationalized triaxial signatures from
VBL-VA001.}
\label{tab:vblcalibration}
\end{table}

\begin{proposition}[Data-grounded rational certificate]
\label{prop:vblcertificate}
For the five rational nodes and signatures in
Table~\ref{tab:vblcalibration}, there is a real-symmetric quadratic polynomial
$P$ satisfying $P(x_j)v_j=v_j$ exactly and
\[
 [C_k,C_\ell]\ne0\quad\text{for some }k,\ell,
 \qquad
 \sup_{-1\le x\le-929/1000}\|P(x)\|_{\op}<\frac{24}{25}.
\]
Every real-symmetric quadratic whose coefficients have a common nontrivial
invariant subspace and which satisfies the same five calibrations has leakage
at least one.  In the pairwise-commuting subcase it is identically $I_3$.
\end{proposition}

\begin{proof}
All ten $3$ by $3$ target determinants are nonzero in exact rational
arithmetic; the smallest normalized absolute determinant is
$0.009489456048\ldots$.  Lemma~\ref{lem:commuting}, with $d=3$, $N=2$ and
$M=5$, therefore proves the common-invariant-line conclusion.  Any common
nontrivial invariant subspace of real-symmetric $3$ by $3$ coefficients yields
such a line by orthogonal complementation, proving the stated reducible-class
bound and the stronger identity conclusion for commuting coefficients.

For the upper bound, fix the $(1,1)$, $(1,2)$ and $(2,2)$ entries of $C_2$ to
$-0.316171$, $-0.537470$ and $-0.909040$, respectively, and solve the remaining
$15$ rational linear equations for the other symmetric entries.  The exact
replay verifies the calibrations and a nonzero commutator.  On the $101$-point
rational grid of $[-1,-929/1000]$, Sylvester's criterion gives
\[
  (191/200)^2I_3-P(x_i)^2\succ0.
\]
The same replay proves
$\sup_{|x|\le1}\|P'(x)\|_{\op}\le5.740500963$.  The grid spacing is
$71/100000$, hence every point is within half a grid spacing of some $x_i$ and
\[
 \|P(x)\|_{\op}
 <\frac{191}{200}
   +5.740500963\frac{71}{200000}
 <0.957038<\frac{24}{25}.
\]
Every comparison in this paragraph is performed on integers and rational
fractions; the printed decimals only abbreviate upward-valid bounds.
\end{proof}

The held-out signatures lie between $0.129^\circ$ and $0.460^\circ$ from their
training counterparts.  Without refitting, the rational witness has maximum
held-out residual $0.0038491$ and dense-grid leakage $0.94870$.  The numerical
comparators in Table~\ref{tab:vblresults} use the unrounded measured nodes and
signatures.  A general nonsymmetric MIMO polynomial attains much lower training
leakage, but its held-out residual is $0.1045$ and its largest coefficient
Frobenius norm is $36.06$, compared with $2.623$ for the rational symmetric
witness.  This single split suggests a reciprocity--robustness tradeoff; it
does not establish a statistical generalization theorem.

For a task-level replay, we froze the full $3$ by $3$ training and held-out
cospectral matrices at the five peaks, applied $P(x_j)G_jP(x_j)$ on the
held-out split, and recomputed the FDD leading eigenpair.  Across the five
frequencies, the angle between the unfiltered and filtered leading directions
is at most $0.2217^\circ$; the leading eigenvalue ratio differs from one by at
most $2.40\times10^{-5}$.  These figures are outputs of the offline replay,
not consequences of the exact certificate.  Lemma~\ref{lem:fddinvariance}
explains the specification: exact calibration preserves an ideal modal
component, while the held-out calculation measures the effect of residual
components and train--test drift.

\begin{table}[H]
\centering
\small
\begin{tabular}{@{}lrrr@{}}
\toprule
Filter & Stop leakage & Fit tolerance & Held-out residual\\
\midrule
Rational symmetric witness & $<0.96000$ & exact & $0.003849$\\
Symmetric SDP optimum & $0.94867$ & exact & $0.003849$\\
General nonsymmetric SDP & $0.14514$ & exact & $0.10447$\\
Scalar quadratic optimum & $0.97520$ & $0.01$ & $0.01000$\\
Best of 64 fixed bases & $0.96966$ & $0.01$ & $0.01011$\\
Exact reducible symmetric class & $\ge1$ & exact & $0$\\
\bottomrule
\end{tabular}
\caption{Measured-calibration controls.  The 64-basis row is a seeded stress
test, not a global optimization certificate over orthogonal bases.}
\label{tab:vblresults}
\end{table}

Thus the measured instance closes both the purely constructed-data objection
and the objection that tangential preservation has no task meaning: its
calibration and validation use disjoint sensor records, FDD supplies the
mode-shape objective, and the exact $0.96$-versus-$1$ gap is not
solver-dependent.  It does not show improved fault diagnosis, damping
estimation, hardware cost or deployment performance; those claims are outside
the present experiment.

\subsection{A failed graph-denoising application gate}

We also tested the standard MIMO graph-filter form
$Y=\sum_{k=0}^{2}S^kXC_k$ on two smooth fields over an $18$ by $23$ grid
graph, using four low-frequency full-spark row signatures.  The symmetric MIMO
SDP collapsed to the identity (validated stop norm $0.99999992$), while the
one-point-normalized scalar Chebyshev filter had stop norm $1/17$.  After two
iterations its best globally rescaled pass error was only $0.00524$ and its
desired-subspace sine was $0.0292$, compared with $0.997$ for the symmetric
MIMO solution.  A nonsymmetric MIMO fit eventually preserved the exact
geometry, but was inferior at short iteration counts.  Thus ordinary graph
denoising does not justify the exact tangential specification, and no such
application claim is made.  We retain this negative result to prevent selecting
an example merely because it matches the theorem's notation.

\section{Relation to established work and scope of the claim}

Tangential matrix-polynomial interpolation is classical
\cite{BallKang1990,Fuhrmann2010}.  Structured matrix-polynomial interpolation
with Hermitian, positive-semidefinite and related symmetries is also established
\cite{AlpayLewkowicz2014}; that work treats full matrix values and lists the
tangential structured variant as a future problem.  Matrix-valued
Nevanlinna--Pick theory handles
supremum-norm analytic interpolation, including complexity constraints
\cite{BlomqvistEtAl2003,BallBolotnikov2016}.  Kuroiwa and Lindquist treat the
bitangential problem with explicit McMillan-degree constraints
\cite{KuroiwaLindquist2006}.  Most directly,
Stefanovski and Georgijevi\'c prove that real stable rational matrices can meet
bitangential constraints while having arbitrarily small spectral norm on a
proper frequency region \cite{StefanovskiGeorgijevic2016}; their degree grows
and they do not compare commuting and noncommuting Hermitian polynomial
coefficients.  Thus small regional norm under matrix tangential constraints is
not itself our contribution.  Tangential interpolation is also a standard
bridge between MIMO model reduction and Krylov projection
\cite{GallivanEtAl2004}; modern block rational approximation includes
matrix-valued weights, noisy-data comparisons, and tangential interpolation
mechanisms \cite{GoseaGuettel2021}.  Recent iterative MIMO schemes even select
low-rank tangential data using maximum-error criteria \cite{JonasBamieh2026}.
Their maximum-error step selects the next interpolation frequency, whereas
their proved monotonicity concerns a weighted $H_2$ error; monotone $H_\infty$
bounds are stated as future work.  Gallivan et al.'s main result instead gives
a generically unique minimal-McMillan-degree rational interpolant.  Neither
result provides a fixed-degree Hermitian-polynomial minimax separation from the
larger class with a common invariant channel subspace.  Minimax polynomial filtering and SDP
formulations have a long numerical-linear-algebra history
\cite{TohTrefethen1998}, and matrix-polynomial spaces arise naturally in block
Krylov methods \cite{RinelliVandebril2026}.  General MIMO graph filters already
use matrix taps in precisely the form $\sum_k S^kXC_k$
\cite{GamaEtAl2018}, while scalar FIR minimax design by convex optimization is
classical \cite{WuBoydVandenberghe1998}.  Kootsookos already formulates
fixed-length MIMO FIR approximation of matrix transfer functions in the
uniform $H_\infty$ norm, proves global error lower bounds, and gives MIMO and
linear-phase design algorithms \cite{Kootsookos1991}.  More closely, Ljungars
and Fu optimize
the operator-norm error of matrix-valued multi-channel linear-phase FIR filters
on a frequency grid by SDP \cite{LjungarsFu1998}.  Their formulation has neither
directional pass interpolation nor an invariant-subspace/reducibility
comparator.  Accordingly, fixed-length MIMO $H_\infty$ approximation, MIMO
linear phase, operator-norm minimax design, and SDP are not claimed as new.

The candidate novelty is therefore narrower than norm-constrained or
degree-constrained tangential interpolation or FIR design alone.  Under fixed
ordinary degree---equivalently fixed reciprocal FIR latency---a three-channel
Hermitian-coefficient tangential minimax problem exhibits an exponential
separation between irreducible coefficients and every reducible symmetric
family.  The excluded family is strictly larger than fixed-basis scalar banks:
it permits a frequency-dependent, noncommuting $2$ by $2$ block, provided one
channel line remains invariant.  The quantitative result also identifies a
joint calibration/reducibility scale: an exponentially small sum of tangential
residual and sampled off-block coupling is sufficient for the constructed
separation, whereas a universal scalar construction proves that an
exponentially small calibration residual can erase every such separation.  We
have not found this two-sided law or the reducible/irreducible separation in
the sources above.

The complete texts of two close restricted papers---Fuhrmann (2010) and
Stefanovski--Georgijevi\'c (2016)---could not be obtained.  Publisher portals,
repository copies, metadata APIs, and independent network routes were checked.
Their accessible statements already rule out broader novelty claims, but
absence of a theorem from inaccessible text is not evidence of priority.  The
claim is therefore limited to the displayed theorem and to the literature
actually inspected.
The FDD lemma and held-out replay support a preprocessing interpretation; they
do not constitute a deployment or diagnostic study.

\paragraph{Reproducibility.}
The source and research ledger are in the
\href{https://github.com/lluiseriksson/finite-sample-spectral-certificates/tree/research/endogenous-visibility}{project repository}.
The repository links the
\href{https://github.com/lluiseriksson/finite-sample-spectral-certificates/blob/research/endogenous-visibility/verification/verify_constant_channel_scaling.py}{three-channel theorem arithmetic replay},
the
\href{https://github.com/lluiseriksson/finite-sample-spectral-certificates/blob/research/endogenous-visibility/verification/verify_quadratic_filter_witness.py}{exact rational replay},
the
\href{https://github.com/lluiseriksson/finite-sample-spectral-certificates/blob/research/endogenous-visibility/research/route_c_block_krylov_benchmark.py}{block-iteration benchmark},
its
\href{https://github.com/lluiseriksson/finite-sample-spectral-certificates/blob/research/endogenous-visibility/results/route_c/block_krylov_benchmark.json}{machine-readable output},
the deliberately negative
\href{https://github.com/lluiseriksson/finite-sample-spectral-certificates/blob/research/endogenous-visibility/results/route_c/graph_filter_pilot.json}{graph-filter pilot},
the
\href{https://github.com/lluiseriksson/finite-sample-spectral-certificates/blob/research/endogenous-visibility/data/route_c/vbl_va001_calibration.json}{VBL-VA001 provenance manifest},
its
\href{https://github.com/lluiseriksson/finite-sample-spectral-certificates/blob/research/endogenous-visibility/verification/verify_vbl_va001_witness.py}{exact continuum replay},
the
\href{https://github.com/lluiseriksson/finite-sample-spectral-certificates/blob/research/endogenous-visibility/research/route_c_vibration_calibration.py}{offline baseline replay},
and its
\href{https://github.com/lluiseriksson/finite-sample-spectral-certificates/blob/research/endogenous-visibility/results/route_c/vbl_va001_calibration.json}{machine-readable output},
the
\href{https://github.com/lluiseriksson/finite-sample-spectral-certificates/blob/research/endogenous-visibility/programme/ROUTE_C_PRIORITY_MATRIX.md}{clause-by-clause priority matrix},
the
\href{https://github.com/lluiseriksson/finite-sample-spectral-certificates/blob/research/endogenous-visibility/formal/SpectralCertificates/RootCount.lean}{Lean root-count proof},
and the
\href{https://github.com/lluiseriksson/finite-sample-spectral-certificates/blob/research/endogenous-visibility/programme/ARTIFACT_MANIFEST.md}{artifact manifest with the PDF hash}.

\footnotesize
\bibliographystyle{plain}
\bibliography{references}
\end{document}
